% chap02_containers.tex - Lists, dicts, sets, and tuples
% ============================================================================

\chapter{Container Types}
\label{chap:containers}

So far, each variable holds one value. But scientific data comes in groups: a series of measurements, a collection of samples, a table of properties. Python's container types hold multiple values in a single variable. This chapter covers the four main ones: lists, tuples, dictionaries, and sets.

% ----------------------------------------------------------------------------
\section{Lists}
\label{sec:lists}

A list is an ordered collection of items. Create one with square brackets:

\begin{lstlisting}
>>> temperatures = [20.1, 21.3, 19.8, 22.0, 20.5]
>>> temperatures
[20.1, 21.3, 19.8, 22.0, 20.5]
>>> elements = ["H", "He", "Li", "Be", "B"]
>>> elements
['H', 'He', 'Li', 'Be', 'B']
\end{lstlisting}

Lists can contain any type of data, and items don't have to be the same type (though they usually are):

\begin{lstlisting}
>>> mixed = [1, "hello", 3.14, True]
>>> mixed
[1, 'hello', 3.14, True]
\end{lstlisting}

\subsection{Indexing}

Access individual items by position using square brackets. The position number is called the \emph{index}. Python uses zero-based indexing---the first item is at index 0, not index 1. This is a common convention in programming (though not universal), so get used to it:

\begin{lstlisting}
>>> elements = ["H", "He", "Li", "Be", "B"]
>>> elements[0]       # first item
'H'
>>> elements[1]       # second item
'He'
>>> elements[4]       # fifth item (last in this list)
'B'
>>> elements[-1]      # last item (negative indexing)
'B'
>>> elements[-2]      # second to last
'Be'
\end{lstlisting}

\begin{warning}
Accessing an index that doesn't exist causes an error:
\begin{lstlisting}
>>> elements[10]
IndexError: list index out of range
\end{lstlisting}
Always check your list length if you're unsure.
\end{warning}

\subsection{Slicing}

Extract a sublist using slice notation \py{[start:stop]}:

\begin{lstlisting}
>>> numbers = [0, 1, 2, 3, 4, 5, 6, 7, 8, 9]
>>> numbers[2:5]      # items from index 2 up to (not including) 5
[2, 3, 4]
>>> numbers[:4]       # from beginning to index 4
[0, 1, 2, 3]
>>> numbers[6:]       # from index 6 to end
[6, 7, 8, 9]
>>> numbers[::2]      # every second item (step of 2)
[0, 2, 4, 6, 8]
>>> numbers[::-1]     # reversed
[9, 8, 7, 6, 5, 4, 3, 2, 1, 0]
\end{lstlisting}

The pattern is \py{[start:stop:step]}. Any part can be omitted:
\begin{itemize}
    \item Omit \py{start}: Defaults to begin at 0
    \item Omit \py{stop}: Defaults to go to the end
    \item Omit \py{step}: Defaults to use step of 1
\end{itemize}

\begin{note}
The \py{stop} index is exclusive---the slice includes items up to but not including that index. This feels odd at first but makes arithmetic cleaner: \py{numbers[0:5]} contains exactly 5 items.
\end{note}

\begin{ponder}
If we have some list \py{numbers}, we get every value back if we use \py{numbers[0:-1:1]}. Why does \py{numbers[::-1]} reverse this list?
\end{ponder}

\subsection{List Length}

The \py{len()} function returns the number of items:

\begin{lstlisting}
>>> elements = ["H", "He", "Li", "Be", "B"]
>>> len(elements)
5
>>> len([])           # empty list
0
\end{lstlisting}

\subsection{Modifying Lists}

Lists are \emph{mutable}, meaning you can change them after creation. (``Mutable'' comes from the same root as ``mutation''---things that can change.) You can modify individual items, add new ones, or remove existing ones.

\textbf{Change an item:}
\begin{lstlisting}
>>> temps = [20.0, 21.0, 22.0]
>>> temps[1] = 21.5
>>> temps
[20.0, 21.5, 22.0]
\end{lstlisting}

\textbf{Add items:}

Here we'll use \emph{methods}---functions that belong to a specific value. You call them using a dot: \py{temps.append(23.0)} means ``call the append function that belongs to temps.'' Methods are actions that the value knows how to perform on itself.

\begin{lstlisting}
>>> temps.append(23.0)        # add to end
>>> temps
[20.0, 21.5, 22.0, 23.0]
>>> temps.insert(0, 19.0)     # insert at position 0
>>> temps
[19.0, 20.0, 21.5, 22.0, 23.0]
>>> temps.extend([24.0, 25.0])  # add multiple items
>>> temps
[19.0, 20.0, 21.5, 22.0, 23.0, 24.0, 25.0]
\end{lstlisting}

\textbf{Remove items:}
\begin{lstlisting}
>>> temps.pop()               # remove and return last item
25.0
>>> temps
[19.0, 20.0, 21.5, 22.0, 23.0, 24.0]
>>> temps.pop(0)              # remove and return item at index 0
19.0
>>> temps
[20.0, 21.5, 22.0, 23.0, 24.0]
>>> temps.remove(21.5)        # remove first occurrence of value
>>> temps
[20.0, 22.0, 23.0, 24.0]
\end{lstlisting}

\begin{warning}
\py{append()} and \py{extend()} look similar but differ:
\begin{lstlisting}
>>> a = [1, 2, 3]
>>> a.append([4, 5])      # adds the list as one item
>>> a
[1, 2, 3, [4, 5]]
>>> b = [1, 2, 3]
>>> b.extend([4, 5])      # adds each item individually
>>> b
[1, 2, 3, 4, 5]
\end{lstlisting}
\end{warning}

\subsection{List Operations}

\begin{lstlisting}
>>> [1, 2, 3] + [4, 5, 6]     # concatenation
[1, 2, 3, 4, 5, 6]
>>> [1, 2] * 3                 # repetition
[1, 2, 1, 2, 1, 2]
>>> 3 in [1, 2, 3, 4, 5]      # membership test
True
>>> 10 in [1, 2, 3, 4, 5]
False
\end{lstlisting}

\subsection{Useful List Methods}

\begin{lstlisting}
>>> numbers = [3, 1, 4, 1, 5, 9, 2, 6]
>>> numbers.count(1)          # how many 1s?
2
>>> numbers.index(5)          # position of first 5
4
>>> numbers.sort()            # sort in place
>>> numbers
[1, 1, 2, 3, 4, 5, 6, 9]
>>> numbers.reverse()         # reverse in place
>>> numbers
[9, 6, 5, 4, 3, 2, 1, 1]
\end{lstlisting}

For a sorted copy without modifying the original, use \py{sorted()}:

\begin{lstlisting}
>>> original = [3, 1, 4, 1, 5]
>>> sorted_copy = sorted(original)
>>> sorted_copy
[1, 1, 3, 4, 5]
>>> original                  # unchanged
[3, 1, 4, 1, 5]
\end{lstlisting}

\begin{labexample}
Processing a list of absorbance measurements:
\begin{lstlisting}
>>> absorbances = [0.452, 0.487, 0.523, 0.478, 0.501]
>>> print(f"Number of measurements: {len(absorbances)}")
Number of measurements: 5
>>> print(f"First reading: {absorbances[0]}")
First reading: 0.452
>>> print(f"Last reading: {absorbances[-1]}")
Last reading: 0.501
>>> print(f"Min: {min(absorbances)}, Max: {max(absorbances)}")
Min: 0.452, Max: 0.523
>>> print(f"Average: {sum(absorbances) / len(absorbances):.3f}")
Average: 0.488
\end{lstlisting}
The built-in functions \py{min()}, \py{max()}, and \py{sum()} work on lists of numbers.
\end{labexample}

% ----------------------------------------------------------------------------
\section{Tuples}
\label{sec:tuples}

A tuple is like a list, but \emph{immutable}---once created, you can't change it. (``Immutable'' is the opposite of mutable: it cannot be mutated or changed.) This might seem like a limitation, but it's often useful. Create tuples with parentheses:

\begin{lstlisting}
>>> coordinates = (3.5, 2.1, 7.8)
>>> coordinates
(3.5, 2.1, 7.8)
>>> rgb = (255, 128, 0)
>>> rgb
(255, 128, 0)
\end{lstlisting}

Access items the same way as lists:

\begin{lstlisting}
>>> coordinates[0]
3.5
>>> coordinates[1:]
(2.1, 7.8)
>>> len(coordinates)
3
\end{lstlisting}

But you cannot modify them:

\begin{lstlisting}
>>> coordinates[0] = 4.0
TypeError: 'tuple' object does not support item assignment
\end{lstlisting}

\subsection{Why Use Tuples?}

If tuples are just restricted lists, why bother? Several reasons:

\begin{enumerate}
    \item \textbf{Intent:} A tuple signals that the data shouldn't change. Coordinates of a point, RGB values of a color---these are fixed groupings.
    \item \textbf{Dictionary keys:} Tuples can be dictionary keys; lists cannot (we'll see why soon).
    \item \textbf{Performance:} Tuples are slightly faster and use less memory (rarely matters in practice).
    \item \textbf{Unpacking:} Tuples are often used for returning multiple values from functions.
\end{enumerate}

\subsection{Tuple Packing and Unpacking}

Python lets you assign multiple variables at once:

\begin{lstlisting}
>>> point = (3.5, 2.1, 7.8)
>>> x, y, z = point           # unpacking
>>> x
3.5
>>> y
2.1
>>> z
7.8
\end{lstlisting}

This also works without explicit parentheses:

\begin{lstlisting}
>>> a, b = 1, 2               # tuple packing on right, unpacking on left
>>> a
1
>>> b
2
>>> a, b = b, a               # swap values!
>>> a
2
>>> b
1
\end{lstlisting}

\begin{note}
A single-element tuple needs a trailing comma:
\begin{lstlisting}
>>> not_a_tuple = (42)        # just a number in parentheses
>>> type(not_a_tuple)
<class 'int'>
>>> actual_tuple = (42,)      # trailing comma makes it a tuple
>>> type(actual_tuple)
<class 'tuple'>
\end{lstlisting}
\end{note}

% ----------------------------------------------------------------------------
\section{Dictionaries}
\label{sec:dicts}

A dictionary stores key-value pairs. Instead of accessing items by position, you access them by name:

\begin{lstlisting}
>>> element = {"name": "Carbon", "symbol": "C", "atomic_number": 6}
>>> element["name"]
'Carbon'
>>> element["atomic_number"]
6
\end{lstlisting}

Create dictionaries with curly braces. Each entry has a key, a colon, and a value. Keys are typically strings, but can be any immutable type.

\subsection{Adding and Modifying}

\begin{lstlisting}
>>> element["mass"] = 12.011      # add new key-value pair
>>> element
{'name': 'Carbon', 'symbol': 'C', 'atomic_number': 6, 'mass': 12.011}
>>> element["mass"] = 12.0        # modify existing value
>>> element["mass"]
12.0
\end{lstlisting}

\subsection{Accessing Values Safely}

Accessing a missing key raises an error:

\begin{lstlisting}
>>> element["electrons"]
KeyError: 'electrons'
\end{lstlisting}

Use \py{.get()} to provide a default instead:

\begin{lstlisting}
>>> element.get("electrons")          # returns None if missing
>>> element.get("electrons", "unknown")  # returns "unknown" if missing
'unknown'
>>> element.get("name", "unknown")    # returns actual value if present
'Carbon'
\end{lstlisting}

Check if a key exists:

\begin{lstlisting}
>>> "name" in element
True
>>> "electrons" in element
False
\end{lstlisting}

\subsection{Dictionary Methods}

\begin{lstlisting}
>>> element.keys()            # all keys
dict_keys(['name', 'symbol', 'atomic_number', 'mass'])
>>> element.values()          # all values
dict_values(['Carbon', 'C', 6, 12.0])
>>> element.items()           # key-value pairs as tuples
dict_items([('name', 'Carbon'), ('symbol', 'C'), ...])
\end{lstlisting}

Remove entries:

\begin{lstlisting}
>>> element.pop("mass")       # remove and return value
12.0
>>> element
{'name': 'Carbon', 'symbol': 'C', 'atomic_number': 6}
>>> del element["symbol"]     # remove without returning
>>> element
{'name': 'Carbon', 'atomic_number': 6}
\end{lstlisting}

\subsection{Nested Dictionaries}

Values can be any type, including other dictionaries:

\begin{lstlisting}
>>> periodic_table = {
...     "H": {"name": "Hydrogen", "mass": 1.008, "number": 1},
...     "He": {"name": "Helium", "mass": 4.003, "number": 2},
...     "Li": {"name": "Lithium", "mass": 6.941, "number": 3}
... }
>>> periodic_table["H"]["mass"]
1.008
>>> periodic_table["Li"]["name"]
'Lithium'
\end{lstlisting}

\begin{labexample}
Storing experimental conditions:
\begin{lstlisting}
>>> experiment = {
...     "sample_id": "RXN-2024-001",
...     "temperature_K": 298.15,
...     "pressure_atm": 1.0,
...     "solvent": "water",
...     "concentrations_M": {
...         "NaCl": 0.1,
...         "HCl": 0.01
...     },
...     "notes": "Standard conditions"
... }
>>> print(f"Sample: {experiment['sample_id']}")
Sample: RXN-2024-001
>>> print(f"NaCl conc: {experiment['concentrations_M']['NaCl']} M")
NaCl conc: 0.1 M
\end{lstlisting}
Dictionaries are excellent for structured data where items have meaningful names.
\end{labexample}

% ----------------------------------------------------------------------------
\section{Sets}
\label{sec:sets}

A set is an unordered collection of unique elements. Sets automatically remove duplicates:

\begin{lstlisting}
>>> elements = {"H", "C", "N", "O", "C", "H"}    # note duplicates
>>> elements
{'O', 'C', 'N', 'H'}                             # duplicates removed
>>> len(elements)
4
\end{lstlisting}

The order may vary---sets are unordered.

\subsection{Set Operations}

Sets support mathematical set operations:

\begin{lstlisting}
>>> group1 = {"H", "C", "N", "O"}
>>> group2 = {"C", "O", "S", "P"}
>>> group1 | group2               # union (all elements in either)
{'S', 'N', 'C', 'O', 'P', 'H'}
>>> group1 & group2               # intersection (elements in both)
{'C', 'O'}
>>> group1 - group2               # difference (in group1 but not group2)
{'N', 'H'}
>>> group1 ^ group2               # symmetric difference (in one but not both)
{'S', 'N', 'P', 'H'}
\end{lstlisting}

\subsection{Membership Testing}

Sets are optimized for membership testing---checking if something is in the set:

\begin{lstlisting}
>>> elements = {"H", "C", "N", "O", "S", "P"}
>>> "C" in elements
True
>>> "Fe" in elements
False
\end{lstlisting}

This is much faster than lists for large collections.

\subsection{Modifying Sets}

\begin{lstlisting}
>>> elements = {"H", "C", "N"}
>>> elements.add("O")             # add one element
>>> elements
{'O', 'C', 'N', 'H'}
>>> elements.update(["S", "P"])   # add multiple
>>> elements
{'S', 'O', 'C', 'P', 'N', 'H'}
>>> elements.remove("S")          # remove (error if missing)
>>> elements.discard("Fe")        # remove (no error if missing)
\end{lstlisting}

\begin{note}
To create an empty set, use \py{set()}, not \py{\{\}}:
\begin{lstlisting}
>>> empty_dict = {}
>>> type(empty_dict)
<class 'dict'>
>>> empty_set = set()
>>> type(empty_set)
<class 'set'>
\end{lstlisting}
Curly braces without colons create an empty dictionary, not a set.
\end{note}

\begin{labexample}
Finding common elements between two samples:
\begin{lstlisting}
>>> sample_a = {"Fe", "Cu", "Zn", "Pb", "Cd"}
>>> sample_b = {"Cu", "Ag", "Au", "Pb"}
>>> common = sample_a & sample_b
>>> print(f"Elements in both samples: {common}")
Elements in both samples: {'Cu', 'Pb'}
>>> only_in_a = sample_a - sample_b
>>> print(f"Only in sample A: {only_in_a}")
Only in sample A: {'Fe', 'Zn', 'Cd'}
\end{lstlisting}
\end{labexample}

% ----------------------------------------------------------------------------
\section{Choosing the Right Container}
\label{sec:choosing-container}

With four container types, how do you choose?

\begin{table}[h]
\centering
\begin{tabular}{lcccc}
\toprule
Characteristic & List & Tuple & Dict & Set \\
\midrule
Ordered & Yes & Yes & No* & No \\
Mutable & Yes & No & Yes & Yes \\
Duplicates allowed & Yes & Yes & Keys: No & No \\
Access by index & Yes & Yes & No & No \\
Access by key & No & No & Yes & No \\
\bottomrule
\end{tabular}
\caption{Container type comparison. *Dicts preserve insertion order since Python 3.7, but aren't indexed by position.}
\label{tab:container-comparison}
\end{table}

\textbf{Use a list when:}
\begin{itemize}
    \item Order matters
    \item You need to add, remove, or change items
    \item Items are accessed by position
    \item Example: time series of measurements
\end{itemize}

\textbf{Use a tuple when:}
\begin{itemize}
    \item Order matters but data shouldn't change
    \item You're grouping related fixed values
    \item Example: (x, y, z) coordinates, RGB colors
\end{itemize}

\textbf{Use a dictionary when:}
\begin{itemize}
    \item Items have meaningful names
    \item You need to look up values by key
    \item Data is structured with labeled fields
    \item Example: element properties, experiment metadata
\end{itemize}

\textbf{Use a set when:}
\begin{itemize}
    \item You need unique elements only
    \item Order doesn't matter
    \item You're doing membership tests or set operations
    \item Example: elements present in a sample, unique identifiers
\end{itemize}

% ----------------------------------------------------------------------------
\section{Converting Between Types}
\label{sec:container-conversion}

Convert between container types using their constructors:

\begin{lstlisting}
>>> list((1, 2, 3))           # tuple to list
[1, 2, 3]
>>> tuple([1, 2, 3])          # list to tuple
(1, 2, 3)
>>> set([1, 2, 2, 3, 3, 3])   # list to set (removes duplicates)
{1, 2, 3}
>>> list({1, 2, 3})           # set to list
[1, 2, 3]
>>> list({"a": 1, "b": 2})    # dict to list (gets keys)
['a', 'b']
>>> list({"a": 1, "b": 2}.values())  # dict values to list
[1, 2]
\end{lstlisting}

A common pattern: remove duplicates from a list while preserving some order:

\begin{lstlisting}
>>> items = [3, 1, 4, 1, 5, 9, 2, 6, 5, 3, 5]
>>> unique = list(dict.fromkeys(items))  # preserves first occurrence order
>>> unique
[3, 1, 4, 5, 9, 2, 6]
\end{lstlisting}

% ----------------------------------------------------------------------------
\section{Nested Structures}
\label{sec:nested}

Containers can hold other containers, creating complex data structures:

\begin{lstlisting}
>>> # List of dictionaries (common for tabular data)
>>> samples = [
...     {"id": "A1", "concentration": 0.1, "absorbance": 0.23},
...     {"id": "A2", "concentration": 0.2, "absorbance": 0.45},
...     {"id": "A3", "concentration": 0.3, "absorbance": 0.68},
... ]
>>> samples[1]["absorbance"]
0.45

>>> # Dictionary with list values
>>> spectra = {
...     "wavelengths": [400, 450, 500, 550, 600],
...     "absorbances": [0.12, 0.34, 0.56, 0.43, 0.21]
... }
>>> spectra["wavelengths"][2]
500
\end{lstlisting}

\begin{labexample}
A complete experiment record:
\begin{lstlisting}
>>> kinetics_run = {
...     "reaction": "A + B -> C",
...     "temperature_K": 298,
...     "initial_concentrations": {"A": 0.1, "B": 0.2},
...     "time_points_s": [0, 60, 120, 180, 240, 300],
...     "concentration_A": [0.1, 0.082, 0.067, 0.055, 0.045, 0.037],
...     "replicates": 3,
...     "notes": ["Used fresh reagents", "Room temp stable"]
... }
>>> # Access nested data
>>> initial_A = kinetics_run["initial_concentrations"]["A"]
>>> final_A = kinetics_run["concentration_A"][-1]
>>> print(f"[A] dropped from {initial_A} to {final_A} M")
[A] dropped from 0.1 to 0.037 M
\end{lstlisting}
\end{labexample}

% ----------------------------------------------------------------------------
\section{Chapter Summary}
\label{sec:containers-summary}

You've learned Python's four main container types:

\begin{itemize}
    \item \textbf{Lists:} ordered, mutable sequences accessed by index
    \item \textbf{Tuples:} ordered, immutable sequences for fixed groupings
    \item \textbf{Dictionaries:} key-value pairs for structured data
    \item \textbf{Sets:} unordered collections of unique elements
\end{itemize}

Key operations:

\begin{itemize}
    \item \textbf{Indexing:} \py{items[0]}, \py{items[-1]}
    \item \textbf{Slicing:} \py{items[1:4]}, \py{items[::2]}
    \item \textbf{Length:} \py{len(items)}
    \item \textbf{Membership:} \py{x in items}
    \item \textbf{List methods:} \py{append()}, \py{extend()}, \py{pop()}, \py{sort()}
    \item \textbf{Dict access:} \py{d["key"]}, \py{d.get("key", default)}
    \item \textbf{Set operations:} \py{|}, \py{\&}, \py{-} (union, intersection, difference)
\end{itemize}

\begin{ponder}
\begin{enumerate}
    \item Why can't lists be dictionary keys, but tuples can?
    \item You have a list of 1000 sample IDs and need to check if a new ID already exists. Would a list or set be faster for this? Why?
    \item When would you use a list of dictionaries versus a dictionary of lists?
    \item What happens if you try to access \py{my\_list[len(my\_list)]}? Why?
\end{enumerate}
\end{ponder}
